\chapter{Parallel Worldline Numerics in \texorpdfstring{\acs{CUDA}}{CUDA}} 

\label{ch:WLNumerics}

\acresetall

\begin{summary}
	In this chapter, I give an overview of the \ac{WLN} 
	technique, and discuss the parallel 
	\ac{CUDA} implementation of the algorithm created for this 
	thesis. In the \ac{WLN} technique, we wish to 
	generate an ensemble of representative closed-loop particle 
	trajectories, and use these to compute an approximate 
	average value for Wilson loops. We show how this can be done with 
	a specific emphasis on cylindrically symmetric magnetic fields.
	The fine-grained, massive parallelism provided by 
	the \ac{GPU} architecture results in considerable speedup in 
	computing Wilson loop averages.
\end{summary}

In this chapter, I will give an overview of a numerical technique 
which can be used to compute the effective actions of external 
field configurations.
The technique, called either \acl{WLN} or the Loop Cloud Method,
was first used by Gies and Langfeld~\cite{Gies:2001zp} and has since been
applied to computation of effective actions 
\cite{Gies:2001tj,Langfeld:2002vy,Gies:2005sb,Gies:2005ym,Dunne:2009zz}
and Casimir energies
\cite{Moyaerts:2003ts,Gies:2003cv,PhysRevLett.96.220401}. More recently, 
the technique has also been applied to pair production~\cite{2005PhRvD..72f5001G}
and the vacuum polarization tensor~\cite{PhysRevD.84.065035}.
\ac{WLN} is able to compute quantum effective actions in
the one-loop approximation to all orders in both the coupling and in the
external field, so it is well suited to studying non-perturbative aspects
of \acp{QFT} in strong background fields. Moreover, the technique maintains
gauge invariance and Lorentz invariance.  The key idea of the technique is
that a path integral is approximated as the average of a finite set of $N_l$
representative closed paths (loops) through spacetime. We use a standard
Monte-Carlo procedure to generate loops which have large contributions to
the loop average.


\section{\texorpdfstring{\acs{QED}}{QED} Effective Action on the Worldline}

Worldline numerics is built on the worldline formalism which was initially 
invented by Feynman~\cite{PhysRev.80.440, PhysRev.84.108}. 
Much of the recent interest in this formalism is based on the 
work of Bern and Kosower, who derived it from the infinite string-tension limit of 
string theory and demonstrated that it provided an efficient means for computing 
amplitudes in \ac{QCD}~\cite{PhysRevLett.66.1669}. 
For this reason, the worldline formalism is often referred to as 
`string inspired'. However, the formalism can also be obtained straight-forwardly 
from first-quantized field theory~\cite{1992NuPhB.385..145S}, which is the approach 
we will adopt here.  In this formalism
the degrees of freedom of the field are represented in terms of one-dimensional 
path integrals over an ensemble of closed trajectories.  


We begin with the \ac{QED} effective action expressed in the proper-time 
formalism, (\ref{eqn:proptimelag}),

\ba
	\label{eqn:trln}
	{\rm Tr~ln}\left[\frac{\slashed{p}+e\slashed{A}_\mu^0
	-m}{\slashed{p}-m}\right] &=& -\frac{1}{2}\int d^4x \int_0^\infty
	\frac{dT}{T}e^{-iTm^2} \nonumber \\
	& &\times {\rm tr}\biggr( \bra{x}e^{iT(\slashed{p}
	+e\slashed{A}^0_\mu)^2}\ket{x} 
	- \bra{x}e^{iTp^2}\ket{x}\biggr).
\ea



To evaluate $\bra{x}e^{iT(\slashed{p}_\mu + e\slashed{A}_\mu)^2}\ket{x}$, we recognize that it is simply 
the propagation amplitude $\braket{x,T}{x,0}$ from ordinary quantum mechanics 
with $(\slashed{p}_\mu + e\slashed{A}_\mu)^2$ playing the role of the Hamiltonian.  
We therefore express this factor in its path integral form:

\ba
	\bra{x}e^{iT(\slashed{p}_\mu + e\slashed{A}_mu)^2}\ket{x} &=&  \mathcal{N}
	\int \mathcal{D}x_\rho(\tau) e^{-\int_0^T d\tau \left[\frac{\dot{x}^2(\tau)}{4} 
	+ i A_\rho x^\rho(\tau)\right]} \nonumber \\
	& & \times \Paulispin.
\ea

$\mathcal{N}$ is a normalization constant which we can fix by using our renormalization 
condition that the fermion determinant should vanish at zero external field:

\ba
	\bra{x}e^{iTp^2}\ket{x} & = & \mathcal{N}\int \mathcal{D} 
		x_p(\tau)e^{-\int_0^T d\tau\frac{\dot{x}^2(\tau)}{4}} \\
	&=& \int \frac{d^4p}{(2\pi T)^4}\bra{x}e^{iTp^2}\ket{p}\braket{p}{x} \\
	& = & \frac{1}{(4\pi T)^2},	
\ea

\be
	\mathcal{N}\int \mathcal{D}x_\rho(\tau) e^{-\int_0^T d\tau \frac{\dot{x}^2(\tau)}{4}} 
	= \frac{1}{(4\pi T)^2}.
\ee

We may now write
\ba
	\mathcal{N}\int \mathcal{D}x_\rho(\tau) e^{-\int_0^T d\tau[\frac{\dot{x}^2(\tau)}{4}+iA_\rho x^\rho(\tau)]}
	\Paulispin \nonumber \\
	= \frac{1}{(4\pi T)^2} \left\langle e^{-i\int_0^T d\tau A_\rho x^\rho(\tau)}\Paulispin\right\rangle_x ,
\ea
where
\be
	\label{eqn:meandef}
	\mean{\hat{\mathcal{O}}}_x = \frac{\int \mathcal{D}x_\rho(\tau) \hat{\mathcal{O}} 
	e^{-\int_0^T d\tau\frac{\dot{x}^2(\tau)}{4}}}{\int \mathcal{D}x_\rho(\tau)  
	e^{-\int_0^T d\tau\frac{\dot{x}^2(\tau)}{4}}}
\ee
is the weighted average of the operator $\hat{\mathcal{O}}$ over an ensemble of closed particle 
loops with a Gaussian velocity distribution.

Finally, combining all of the equations from this section results in the 
one-loop effective action for \ac{QED} on the worldline:

\ba
	\label{eqn:QEDWL}
	\Gamma^{(1)}[A_\mu] &=& \frac{2}{(4\pi)^2}\int_0^\infty
	\frac{dT}{T^3}e^{-m^2T}\int d^4x_{\rm CM} \nonumber \\
	& & \times \left[\left\langle e^{i\int_0^Td\tau A_\rho(x_{\rm CM}+x(\tau))\dot{x}^\rho(\tau)}
	\Paulispin\right\rangle _x -1\right]. \nonumber \\
	& &
\ea

\section{Worldline Numerics}

The averages, $\mean{\hat{\mathcal{O}}}$, defined by equation (\ref{eqn:meandef})
involve 
functional integration over every possible closed path through spacetime
which has a Gaussian velocity distribution.  
The prescription of the \ac{WLN} technique is to compute 
these averages approximately using a finite set of $N_l$ representative loops 
on a computer.  The worldline average is then approximated as the mean of 
an operator evaluated along each of the worldlines in the ensemble:

\be
	\mean{\hat{\mathcal{O}}[x(\tau)]} \approx 
	\frac{1}{N_l} \sum_{i=1}^{N_l} \hat{\mathcal{O}}[x_i(\tau)].
\ee


\subsection{Loop Generation} 
\label{sec:loopgen}
The velocity distribution for the loops depends on
the proper time, $T$.  However, generating a separate ensemble of loops for 
each value of $T$ would be very computationally expensive.  This problem is alleviated by generating
a single ensemble of loops, $\vec{y}(\tau)$, representing unit proper time,
and scaling those loops accordingly for different values of $T$:

\be 
\vec{x}(\tau) = \sqrt{T}\vec{y}(\tau/T) ,
\ee

\be \int_0^T d\tau \vec{\dot{x}}^2(\tau) \rightarrow \int_0^1 dt
\vec{\dot{y}}^2(t).
\ee

There is no way to treat the integrals as continuous as we generate our loop
ensembles.  Instead, we treat the integrals as sums over discrete points
along the proper-time interval $[0,T]$.  This is fundamentally different
from space-time discretization, however.  Any point on the worldline loop
may exist at any position in space, and $T$ may take on any value.  It is
important to note this distinction because the Worldline method retains
Lorentz invariance while lattice techniques, in general, do not.

The challenge of loop cloud generation is in generating a discrete set
of points on a unit loop which obeys the prescribed velocity distribution.
There are a number of different algorithms for achieving this goal which have
been discussed in the literature.  Four possible algorithms are compared
and contrasted in \cite{Gies:2003cv}.  In this work, we choose a more
recently developed algorithm, dubbed ``d-loops", which was first described
in \cite{Gies:2005sb}. To generate a ``d-loop", the number of points is iteratively 
doubled, placing the new points in a Gaussian distribution between the existing neighbour points.
We quote the algorithm directly:

\begin{quote} 
  \begin{singlespace}
    \begin{itemize} 
	  \item[(1)] Begin with one arbitrary point
		$N_0=1$, $y_{N}$.

	  \item[(2)] Create an $N_1=2$ loop, i.e., add a point $y_{N/2}$ that is
	    distributed in the heat bath of $y_N$ with
	   \begin{equation} 
		 e^{-\frac{N_1}{4} 2 (y_{N/2} -y_{N})^2}. \label{yn2}
	   \end{equation}

      \item[(3)] Iterate this procedure, creating an $N_k=2^k$ppl loop by
        adding $2^{k-1}$ points $y_{{qN}/{N_k}}$, $q=1,3,\dots, N_k-1$ with
        distribution
      \begin{equation} 
	    e^{-\frac{N_k}{4} 2 [y_{qN/N_k} -\frac{1}{2}(y_{(q+1)N/N_k}+
		y_{(q-1)N/N_k})]^2}. \label{ynk}
	  \end{equation}

      \item[(4)] Terminate the procedure if $N_k$ has reached $N_k=N$ for
        unit loops with $N$ ppl.

      \item[(5)] For an ensemble with common center of mass, shift each
        whole loop accordingly.

    \end{itemize}
  \end{singlespace}	
\end{quote}

The above d-loop algorithm was selected since it is simple and about 10\%
faster than previous algorithms, according to its developers, because it
requires fewer algebraic operations.

The generation of the loops is largely independent from the main program.
Because of this, it was simpler to generate the loops using a Matlab 
script. The Matlab function used for generating d-loops using the 
above algorithm can be found in appendix \ref{sec:srcWLgen}.


This function was used to produce text files containing the worldline data for 
ensembles of loops. Then, these text files were read into memory at the 
launch of each calculation. The results of this generation routine can 
be seen in figure \ref{fig:worldlineplot}.

When the \ac{CUDA} kernel is called\footnote{Please see appendix \ref{ch:cudafication} for an overview of \ac{CUDA}.}, 
every thread in every block executes 
the kernel function with its own unique identifier. So, it is best to 
generate worldlines in integer multiples of the number of threads per 
block. The Tesla C1060 device allows up to 512 threads per block.


\begin{figure}
	\centering
		\includegraphics[width=14cm]{images/worldlineplot.eps}
	\caption[Discretization of the worldline loop]
	{A single discrete worldline loop shown at several levels 
	of discretization. The loops form fractal patterns and have a strong 
	parallel with Brownian motion. The colour
	represents the phase of a particle travelling along the loop, and 
	begins at dark blue, progresses in a random walk through yellow, 
	and ends at dark 
	red. The total flux through this particular worldline at $T=1$ and 
	$B=B_k$ is about $0.08 \pi/e$.}
	\label{fig:worldlineplot}
\end{figure}

\section{Cylindrical Worldline Numerics} 
We now consider cylindrically symmetric external magnetic fields.
In this case, we may simplify (\ref{eqn:QEDWL}),

\be \frac{\Gamma^{(1)}_{\rm ferm}}{T
L_z} = \frac{1}{4\pi} \int_0^\infty \rho_{\rm cm}
  \left[ \int_0^\infty \frac{dT}{T^3}e^{-m^2T}\left\{\langle
  W\rangle_{\vec{r}_{\rm cm}} - 1 -\frac{1}{3}(eB_{\rm
  cm}T)^2\right\}\right]d\rho.
  \label{eqn:cylEA}
\ee 

\subsection{Cylindrical Magnetic Fields} 

We have $\vec{B} =
B(\rho)$\bhattext{z} with

\be \label{eqn:BWLN} B(\rho) = \frac{A_\phi(\rho)}{\rho} +
\frac{dA_\phi(\rho)}{d\rho} \ee if we make the gauge choice that $A_0 =
A_\rho = A_z = 0$.

We begin by considering $\Aphi$ in the form

\be \Aphi = \frac{F}{2\pi \rho}f_\lambda(\rho) \ee so that \be
B_z(\rho)=\frac{F}{2\pi\rho}\frac{df_\lambda(\rho)}{d\rho} \ee and the total
flux is \be \Phi=F(f_\lambda(L_\rho)-f_\lambda(0)).  \ee It is convenient
to express the flux in units of $\frac{2 \pi}{e}$ and define a dimensionless
quantity

\be 
	\mathcal{F}=\frac{e}{2 \pi} F.
\ee

%If we require that $B_z(\rho)$ remains finite at $\rho \to 0$ and that
% $\frac{dB_z(\rho)}{d\rho}\biggr|_{\rho=0}=0$, then we can choose without
%loss of generality that $f_\lambda(\rho) = C_1\rho^2 + C_2\rho^4 +...$
%for small $\rho$ and $f_\lambda(L_\rho) = 1$.

%\be C_1=\frac{eB_z(0)}{2 \mathcal{F}}\equiv\frac{e B_0}{2\mathcal{F}} \ee

%\subsection{The Classical Action}

%\ba \Gamma_{\rm Cl}& =& -\int d^4x\frac{B^2}{2} \\ \frac{\Gamma_{\rm Cl}}{T
%L_z}  & =& \frac{F^2}{2\pi}
%  \int_0^\infty \left( f'_\lambda(\rho^2)\right)^2 d\rho^2
%\ea

%If we take \be f_\lambda(\rho^2) = \frac{\rho^2}{(\lambda^2+\rho^2)} \ee then,

%\be \frac{\Gamma_{\rm cl}}{TL_z} = \frac{-2\pi (\mathcal{F}/e)^2}{3 \lambda^2}
%=\frac{-F^2}{6\pi \lambda^2} \ee

\subsection{Wilson Loop}

The quantity inside the angled brackets in equation (\ref{eqn:QEDWL}) is a 
gauge invariant observable called a Wilson loop. We note that the proper time
integral provides a natural path ordering for this operator.
The Wilson loop expectation value is

\be
\label{eqn:wilsonloop}
\langle W\rangle_{\vec{r}_{\rm cm}}= \left \langle e^{ie\int_0^T d\tau
\vec{A}(\vec{r}_{{\rm cm}} + \vec{r}(\tau)) \cdot \dot{\vec{r}}}
 \frac{1}{4} {\rm tr}  e^{\frac{e}{2}\int_0^T d\tau
  \sigma_{\mu \nu}F_{\mu \nu}(\vec{r}_{{\rm cm}} + \vec{r}(\tau))}\right
  \rangle_{\vec{r}_{\rm cm}},
\ee
which we look at as a product between a scalar part ($e^{ie\int_0^T d\tau
\vec{A}(\vec{r}_{{\rm cm}} + \vec{r}(\tau)) \cdot \dot{\vec{r}}}$)
and a fermionic part ($\frac{1}{4} {\rm tr}  e^{\frac{e}{2}\int_0^T d\tau
  \sigma_{\mu \nu}F_{\mu \nu}(\vec{r}_{{\rm cm}} + \vec{r}(\tau))}$).
  

  
\subsubsection{Scalar Part}

In a magnetic field, the scalar part is related to the flux through 
the loop, $\Phi_B$, by Stokes theorem:

\ba
 e^{ie\int_0^T d\tau
 \vec{A} \cdot \dot{\vec{r}}} &=&
 e^{ie\oint \vec{A}\cdot d\vec{r}} = e^{ie\int_{\vec{\Sigma}} \vec{\nabla}\times\vec{A} \cdot d\vec{\Sigma}}\\
 & = & e^{ie\int_{\vec{\Sigma}} \vec{B} d\vec{\Sigma}} = e^{ie\Phi_B}.
\ea
So, this factor accounts for the Aharonov-Bohm phase aquired by 
particles in the loop.

The loop discretization results in the following approximation of the
scalar integral:

\be
  \oint \vec{A}(\vec{r})\cdot d\vec{r} =  \sum_{i=1}^{N_{ppl}}
  \int_{\vec{r}^i}^{\vec{r}^{i+1}}\vec{A}(\vec{r})\cdot d\vec{r}.
\ee 
Using a linear parameterization of the positions, the line integrals are

\be 
  \int_{\vec{r}^i}^{\vec{r}^{i+1}}\vec{A}(\vec{r})\cdot d\vec{r} =
  \int_0^1dt \vec{A}(\vec{r}(t))\cdot(\vec{r}^{i+1} - \vec{r}^i).
\ee
Using the same gauge choice outlined above ($\vec{A}=A_\phi \hat{\phi}$),
we may write

\be 
	\vec{A}(\vec{r}(t)) = \frac{\mathcal{F}}{e\rho^2} 
	f_\lambda(\rho^2)(-y,x,0),
\ee 
where we have chosen $f_\lambda(\rho^2)$ to depend on $\rho^2$ instead
of $\rho$ to simplify some expressions and to 
avoid taking many costly square roots in the worldline numerics.
We then have

\be \int_{\vec{r}^i}^{\vec{r}^{i+1}}\vec{A}(\vec{r})\cdot
d\vec{r} = \mathcal{F} (x^iy^{i+1}-y^i x^{i+1})\int_0^1
dt\frac{f_\lambda(\rho_i^2(t))}{\rho_i^2(t)}.  \ee The linear interpolation
in Cartesian coordinates gives

\be 
	\label{eqn:rhoi} \rho_i^2(t) = A_i + 2B_it + C_i t^2,
\ee
where
\ba 
	A_i &=& (x^i)^2 + (y^i)^2 \\ 
	B_i&=&x^i(x^{i+1} - x^i) + y^i(y^{i+1}-y^i)\\ 
	C_i &= &(x^{i+1}-x^i)^2 + (y^{i+1}-y^i)^2. 
	\label{eqn:Ci} 
\ea

In performing the integrals along the straight lines connecting
each discretized loop point, we are in danger of violating gauge invariance.
If these integrals can be performed analytically, than gauge invariance
is preserved exactly.  However, in general, we wish to compute these integrals
numerically.  In this case, gauge invariance is no longer guaranteed, but
can be preserved to any numerical precision that's desired.

\subsubsection{Fermion Part} 

For fermions, the Wilson loop
is modified by a factor,

\ba 
W^{\rm ferm.} &=& \frac{1}{4}{\rm tr}\left(e^{\frac{1}{2}
	e\int_0^T d\tau \sigma_{\mu \nu}F^{\mu \nu}}\right)\\ 
	&=& \frac{1}{4}{\rm tr}\left(e^{\sigma_{x y} 
	e\int_0^T d\tau B\left(x(\tau)\right)}\right) \\	
	&=& \cosh{\left(e \int_0^T d\tau B\left(x(\tau)\right)\right)}\\
	&=& \cosh{\left(2\mathcal{F}\int_0^T
	d\tau f'_\lambda(\rho^2(\tau))\right)},
	\label{eqn:Wfermfpl}
\ea 
where I have used the
relation 
\be 
	eB = 2\mathcal{F}\frac{d f_\lambda(\rho^2)}{d \rho^2} =
	2\mathcal{F}f'_\lambda(\rho^2).  
\ee 

This factor represents an additional contribution to the 
action because of the spin interaction with the magnetic field. 
Classically, for a particle with a magnetic moment $\vec{\mu}$ 
travelling through a magnetic field in a time $T$, the 
action is modified by a term given by
\be
	\Gamma^0_{\rm spin} = \int_0^T \vec{\mu} \cdot \vec{B}(\vec{x}(\tau)) d\tau.
\ee
The magnetic moment is related to the electron spin 
$\vec{\mu} = g\left(\frac{e}{2m}\right)\vec{\sigma}$, 
so we see that the integral in the above quantum fermion factor is 
very closely related to the classical action 
associated with transporting a magnetic moment through a magnetic field:
\be
	\Gamma^0_{\rm spin} = g\left(\frac{e}{2m}\right) \sigma_{x y} \int_0^T B_z(x(\tau))d\tau.
\ee
Qualitatively, we could write
\be 
	W^{\rm ferm} \sim \cosh{\left(\Gamma^0_{\rm spin}\right)}.
\ee


As a possibly useful aside, we may want to express 
the integral in terms of $f_\lambda(\rho^2)$ instead of its derivative.
We can do this by integrating by parts:
\ba 
	\int_0^T d\tau f'_\lambda(\rho^2(\tau)) &=& 
	\frac{T}{N_{\rm ppl}}\sum_{i=1}^{N_{\rm ppl}}\int_0^1 dt f'_\lambda(\rho^2_i(\tau))
	\\ & = & 
	\frac{T}{N_{\rm ppl}}\sum_{i=1}^{N_{\rm ppl}}\biggr[
	\frac{f_\lambda(\rho^2_i(t))}{2(B_i+C_it)}\biggr|^{t=1}_{t=0} \nonumber \\
  & & +\frac{C_i}{2}\int_0^1 \frac{f_\lambda(\rho^2_i(t))}{(B_i+C_i
  t)^2} dt\biggr] \\ & = &  \frac{T}{N_{\rm ppl}}\sum_{i=1}^{N_{\rm
  ppl}}\frac{C_i}{2}\int_0^1 \frac{f_\lambda(\rho^2_i(t))}{(B_i+C_i t)^2} dt,
\ea 
with $\rho_i^2(t)$ given by equations (\ref{eqn:rhoi}) to (\ref{eqn:Ci}).
The second equality is obtained from integration-by-parts.  In the third
equality, we use the loop sum to cancel the boundary terms in pairs:

\be 
\label{eqn:Wfermfl}
W^{{\rm ferm.}} = \cosh{\left(\frac{\mathcal{F}T}{N_{\rm
ppl}}\sum_{i=1}^{N_{{\rm ppl}}}
  C_i \int_0^1 dt \frac{f_\lambda(\rho^2_i(t))}{(B_i+C_i t)^2}\right)}.
\ee
In most cases, one would use equation (\ref{eqn:Wfermfpl}) to compute the fermion factor 
of the Wilson loop. However, 
equation (\ref{eqn:Wfermfl}) may be useful in cases where $f'_\lambda(\rho^2(\tau))$
is not known or is difficult to compute. 

%\subsubsection{Step-Function Flux Tube} Consider the specific example
%of a step-function flux tube \be B_z(\rho) = \frac{F}{\pi \lambda^2}
%\theta(\lambda-\rho) \ee so that

%\be f_\lambda(\rho^2)=\frac{\rho^2}{\lambda^2}\theta(\lambda^2-\rho^2) +
%\theta(\rho^2-\lambda^2).  \ee

\subsection{Renormalization}

The field strength renormalization counter-terms result from the small $T$
behaviour of the worldline integrand.  In the limit where $T$ is very small,
the worldline loops are very localized around their center of mass.  So,
we may approximate their contribution as being that of a constant field
with value $\vec{A}(\vec{r}_{{\rm cm}})$.  Specifically, we require that the field
change slowly on the length scale defined by $\sqrt{T}$. This condition on 
$T$ can be written

\be 
T \ll \left|\frac{m^2}{e B'(\rho^2)}\right| 
= \left| \frac{m^2}{2\mathcal{F}f''_\lambda(\rho^2_{\rm cm})}\right|.
\ee 

When this limit is satisfied, we may use the exact expressions for the
constant field Wilson loops to determine the small $T$ behaviour of the
integrands and the corresponding counter terms.

The Wilson loop averages for constant magnetic fields in scalar and 
fermionic \ac{QED} are

\be
	\mean{W}_{\rm ferm} = eBT\coth{(eBT)}
\ee
and
\be
	\mean{W}_{\rm scal} = \frac{eBT}{\sinh{(eBT)}}.
\ee
So, the integrand for fermionic \ac{QED} in the limit of small $T$ is

\ba \label{eqn:fermI} I_{\rm ferm}(T) &=&
\frac{e^{-m^2T}}{T^3}\left[eB(\vec{r}_{cm})T\coth{(eB(\vec{r}_{cm})T)}
- 1 -\frac{e^2}{3} B^2(\vec{r}_{cm})T^2\right] \nonumber \\ &\approx&-\frac{(eB)^4
T}{45}+\frac{1}{45} (eB)^4 m^2 T^2+\left( \frac{2 (eB)^6}{945}-\frac{(eB)^4
m^4}{90}\right)T^3 \nonumber \\
  & & +\frac{(7 (eB)^4 m^6-4 (eB)^6 m^2)T^4 }{1890}+O(T^5).
\ea
For \ac{ScQED}, we have
\ba 
\label{eqn:scalI} I_{\rm scal}(T) &=&
\frac{e^{-m^2T}}{T^3}\left[eB(\vec{r}_{cm})T/\sinh{(eB(\vec{r}_{cm})T)}
- 1 +\frac{1}{6} (eB)^2(\vec{r}_{cm})T^2\right] \nonumber \\ &\approx&\frac{7
(eB)^4 T}{360}-\frac{7  (eB)^4 m^2 T^2}{360}+\frac{(147 (eB)^4 m^4-31
(eB)^6)T^3 }{15120} \nonumber \\ & &+\frac{ (31 (eB)^6 m^2-49 (eB)^4
m^6)T^4}{15120}+O(T^5). 
\ea

Beyond providing the renormalization conditions, these expansions can be used
in the small $T$ regime to avoid a problem with the Wilson loop uncertainties
in this region.  Consider the uncertainty in the integrand arising from the
uncertainty in the Wilson loop:

\ba 
\delta I(T) &=& \frac{\partial I}{\partial W} \delta W \\
  &=& \frac{e^{-m^2 T}}{T^3} \delta W.
\ea In this case, even though we can compute the Wilson loops for small $T$
precisely, even a small uncertainty is magnified by a divergent factor when
computing the integrand for small values of $T$.  So, in order to perform
the integral, we must replace the small $T$ behaviour of the integrand with
the above expansions (\ref{eqn:fermI}) and (\ref{eqn:scalI}).  Our worldline
integral then proceeds by analytically computing the integral for the small
$T$ expansion up to some small value, $a$, and adding this to the remaining
part of the integral~\cite{MoyaertsLaurent:2004}:

\be 
	\int_0^{\infty} I(T) dT = \underbrace{\int_0^a I(T) dT}_{\rm small ~T} 
	+ \underbrace{\int_a^\infty I(T)dT}_{\rm worldline ~numerics}.
\ee
Because this normalization proceedure uses the constant field expressions for small values of 
$T$, this scheme introduces a small systematic uncertainty. To improve on the 
method outlined here, the derivatives of the background field can 
be accounted for by using the analytic forms of the heat kernel expansion to perform the 
renormalization~\cite{Gies:2001tj}. 

\begin{figure}[ht] 
  \begin{center}
    \includegraphics[width=5in]{images/Integrand.eps} 
	    \caption[small $T$ behaviour of worldline numerics]
		{The small $T$
		behaviour of \ac{WLN}.  The data points represent the numerical
		results, where the error bars are determined from the jackknife analysis described 
		in chapter \ref{ch:WLError}.
		The solid line represents the exact solution while the dotted line represents
		the small $T$ expansion of the exact solution.  Note the amplification of
		the uncertainties.} 
  \end{center} 
\end{figure}

\section{Computing an Effective Action}

The ensemble average in the effective action is simply the sum over the
contributions from each worldline loop, divided by the number of loops in
the ensemble.  Since the computation of each loop is independent of the
other loops, the ensemble average may be straightforwardly parallelized by
generating separate processes to compute the contribution from each loop.
For this parallelization, four Nvidia Tesla C1060 \acp{GPU} were used through
the \ac{CUDA} parallel programing framework.  Because \acp{GPU} can spawn
thousands of parallel processing threads\footnote{Each Tesla C1060 device has 30 multiprocessors
which can each support 1,024 active threads.  So, the number of threads available at a time 
is 30,720 on each of the Tesla devices.  Billions of threads may be scheduled on each 
device~\cite{cudazone}.} 
with much less computational overhead
than an \ac{MPI} cluster, they excel at handling a very large number of parallel 
threads, although the clock speed is slower and fewer memory resources are typically available.
In contrast, parallel computing on a cluster using \acs{CPU} 
tends to have a much higher speed per thread, but there are typically fewer 
threads available.  The worldline technique is exceedingly parallelizable, 
and it is a straightforward matter to divide the task into thousands or tens of thousands of 
parallel threads.  In this case,  one should expect excellent performance 
from a \ac{GPU} over a parallel \ac{CPU} implementation, unless thousands 
of \acp{CPU} are available for the program. The \ac{GPU} architecture has 
recently been used by another group for computing Casimir forces using 
worldline numerics~\cite{2011arXiv1110.5936A}.
Figure \ref{fig:coprocessing} 
illustrates the coprocessing and parallelization scheme used here for the 
worldline numerics.

\begin{figure}
	\centering
		\includegraphics[width=12cm]{images/CoprocessingDiagram.eps}
                \caption[Heterogeneous processing scheme for worldline numerics]
				{The \ac{CPU} manages the loops which compute the 
					integrals over center of mass and proper time. 
					For each proper time integral, we require the results 
					from a large number of individual worldlines. The \ac{GPU}
					is used to compute the integral 
					over each worldline in parallel, and the 
					results are returned to the \ac{CPU} for use in the effective action
					calculation.}.
	\label{fig:coprocessing}
\end{figure}

In an
informal test, a Wilson loop average was computed using an ensemble of 5000
loops in 4.7553 seconds using a serial implementation 
while the \ac{GPU} performed the same calculation
in 0.001315 seconds using a \ac{CUDA} implementation.  
A parallel \ac{CPU} code would require a cluster with thousands of 
cores to achieve a similar speed, even if we assume linear (ideal) 
speed-up.  So, for \ac{WLN} computations,
a relatively inexpensive \ac{GPU} can be expected to outperform a small
or mid-sized cluster.  This increase in computation speed has enabled the
detailed parameter searches discussed in this dissertation.

Of course, there are also tradeoffs from using the 
\ac{GPU} architecture with the \ac{WLN} technique. The most significant 
of these is the limited availability of memory on the device. A \ac{GPU} device 
provides only a few GB of global memory (4GB on the Tesla C1060). This limit 
forces compromises between the number of points-per-loop and the number of 
loops to keep the total size of the loop cloud data small. The limited availability of memory 
resources also limits the number of threads which can be executed concurrently on the device.
Because of the overhead associated with transfering data to and from the device, 
the advantages of a \ac{GPU} over a \ac{CPU} cluster are most pronounced 
on problems which can be divided into several hundred or thousands of processes. 
Therefore, the \ac{GPU} may not offer great performance advantages 
for a small number of loops. Finally, there is some additional complexity involved 
in programming for the \ac{GPU} in terms of learning specialized libraries and 
memory management on the device. This means that the code may take longer to develop
and there may be a learning curve for researchers. However, this problem 
is not much more pronounced with \ac{GPU} programming than with other 
parallelization strategies.

Once the ensemble average of the Wilson loop has been computed, computing
the effective action is a straightforward matter of performing numerical
integrals.  The effective action density is computed by performing the
integration over proper time, $T$.  Then, the effective action is computed
by performing a spacetime integral over the loop ensemble center of mass.
In all cases where a numerical integral was performed, Simpson's method was
used~\cite{burden2001numerical}. Integrals from 0 to $\infty$ were mapped to the 
interval $[0,1]$ using substitutions of the form $x = \frac{1}{1+T/T_{\rm max}}$, 
where $T_{\rm max}$ sets the scale for the peak of the integrand. In the constant field 
case, for the integral over proper time,
 we expect $T_{\rm max} \sim 3/(eB)$ for large fields and $T_{\rm max} \sim 1$ for 
fields of a few times critical or smaller. In chapter \ref{ch:WLError}, we present 
a detailed discussion of how the statistical and discretization uncertainties 
can be computed in this technique.

In this implementation, the numerical integrals are done using a serial CPU
computation. 
This serial portion of the algorithm tends to limit the speedup that can be achieved with 
the large number of parallel threads on the \ac{GPU} device\footnote{By Amdahl's law, 
for a program with a ratio, $P$, of parallel to serial computations on $N$ processors, 
the speedup is given by $S < \frac{1}{(1-P)+\frac{P}{N}}$. This law predicts 
rapidly deminishing returns from increasing the number of processors when $P>0$ 
for a fixed problem size.}. 
However, an important 
benefit of the large number of threads available on the \ac{GPU} is that the 
number of worldlines in the ensemble can be increased without limit, as long 
as more threads are available, without significantly increasing the computation time. 
If perfect occupancy could be achieved on a Tesla C1060 device, an ensemble 
of up to 30,720 worldlines could be computed concurrently.
Thus, the \ac{GPU} provides an excellent architecture for improving the 
statistical uncertainties. Additionally, there is room for further optimization of 
the algorithm by parallelizing the serial portions of the 
algorithm to achieve a greater speedup.

More details about implementing this algorithm on the \ac{CUDA} architecture 
can be found in appendix \ref{ch:cudafication}. A listing of the 
\ac{CUDA} \ac{WLN} code appears in appendix \ref{ch:sourcecode}. 



\section{Verification and Validation}

The \ac{WLN} software can be validated and verified by making sure that it 
produces the correct results where the derivative expansion is a good approximation, 
and that the results are consistent with previous numerical calculations of flux tube 
effective actions. For this reason, the validation was done primarily with flux tubes with 
a profile defined by the function
\be
	f_\lambda(\rho^2) = \frac{\rho^2}{(\lambda^2 + \rho^2)}.
\ee
For large values of $\lambda$, this function varies slowly on the Compton wavelength 
scale, and so the derivative expansion is a good approximation. Also, flux tubes 
with this profile were studied previously using \ac{WLN}
\cite{Moyaerts:2003ts, MoyaertsLaurent:2004}.

Among the results presented in~\cite{MoyaertsLaurent:2004} is a comparison of 
the derivative expansion and \ac{WLN} for this magnetic field 
configuration. The result is that the next-to-leading-order term 
in the derivative expansion is only a small correction to the the 
leading-order term for $\lambda \gg \lambda_e$, where the derivative 
expansion is a good approximation. The derivative expansion breaks
down before it reaches its formal validity limits 
at $\lambda \sim \lambda_e$. For this reason,
we will simply focus on the leading order derivative expansion (which we call the
\ac{LCF} approximation).
The effective action of \ac{QED} in the \ac{LCF} approximation is given 
in cylindrical symmetry by
\ba
	\label{eqn:LCFferm}
	\Gamma^{(1)}_{\rm ferm} &=& \frac{1}{4\pi}\int_0^\infty dT 
	\int_0^\infty \rho_{\rm cm} d\rho_{\rm cm}\frac{e^{-m^2T}}{T^3} \nonumber \\
	& &\left\{eB(\rho_{\rm cm})T\coth{(eB(\rho_{\rm cm})T)} 
	- 1 -\frac{1}{3}(eB(\rho_{\rm cm})T)^2\right\}.
\ea

\begin{figure}
	\centering
		\includegraphics[width=12cm]{images/igrandvsT4.F10l21.eps}
	\caption[Comparison with derivative expansion for $T$ integrand]
	{The integrand of the proper time, $T$, integral for three different values 
	of the radial coordinate, $\rho$ for a $\lambda = 1$ flux tube. The solid lines 
	represent the zeroth-order derivative expansion, which, as expected, is a good approximation 
	until $\rho$ becomes too small.}
	\label{fig:igrandvsT}
\end{figure}

Figure \ref{fig:igrandvsT} shows a comparison between the proper time integrand,
\be
	\frac{e^{-m^2T}}{T^3}\left[\langle W\rangle_{\vec{r}_{\rm cm}} 
		- 1 -\frac{1}{3}(eB_{\rm cm}T)^2\right],
\ee
and the \ac{LCF} approximation result for a flux tube with $\lambda = \lambda_e$ and
$\mathcal{F} = 10$. In this case, the \ac{LCF} approximation is only appropriate far from the 
center of the flux tube, where the field is not changing very rapidly. In the figure, we can 
begin to see the deviation from this approximation, which gets more pronounced closer to the 
center of the flux tube (smaller values of $\rho$).



\begin{figure}
	\centering
		\includegraphics[width=12cm]{images/vsxcm.ferm.smooth.eps}
	\caption[Comparison with \acs{LCF} approximation for action density]
	{The fermion term of the effective action density as a function 
	of the radial coordinate for a flux tube with width $\lambda = 10 \lambda_e$. }
	\label{fig:vsrhocm}
\end{figure}

The effective action density for a slowly varying flux tube is plotted in figure
\ref{fig:vsrhocm} along with the \ac{LCF} approximation. In this case, 
the \ac{LCF} approximation agrees within the statistics of the \ac{WLN}.

%\begin{figure}
%	\centering
%		\includegraphics[width=12cm]{images/EAaf2vsl.eps}
%	\caption[One-loop correction versus flux tube radius]
%	{The fermion term of the effective action for flux tubes 
%	with a wide range of widths and for several values of flux. }
%	\label{fig:EAaf2vsl}
%\end{figure}

%In figure \ref{fig:EAaf2vsl}, we plot the fermionic term of the effective 
%action over several orders of magnitude of $\lambda$ for several values of 
%the dimensionless flux parameter, $\mathcal{F}$. These results are consistent with 
%previous \ac{WLN} calculations~\cite{Moyaerts:2003ts, MoyaertsLaurent:2004}.

\section{Conclusions}

In this chapter, I have reviewed the \ac{WLN} numerical technique with a focus 
on computing the effective action of \ac{QED} in nonhomogeneous, cylindrically symmetric 
magnetic fields. The method uses a Monte Carlo generated ensemble of worldline loops 
to approximate a path integral in the worldline formalism. These worldline loops are 
generated using a simple algorithm and encode the information about the magnetic 
field by computing the flux through the loop and the action aquired from transporting a 
magnetic moment around the loop.
This technique preserves Lorentz symmetry exactly and can preserve 
gauge symmetry up to any required precision. 

I have discussed implementing this technique on \ac{GPU} architecture using 
\ac{CUDA}. The main advantage of this architecture is that it allows for a 
very large number of concurrent threads which can be utilized with 
very little overhead. In practice, this means that a large ensemble 
of worldlines can be computed concurrently, thus allowing for a considerable 
speedup over serial implementations, and allowing for the precision of the 
numerics to increase according to the number of threads available. 
